\section{Detect \VCs}
\label{sec:vc-detection}

This sections details how \tool  
stitches causal relationships 
to form fault propagation chains (\Cref{sec:chain-forming}), 
how to avoid incompatible stitching (\Cref{sec:chain-forming:approx-check}), 
and how \tool searches for cycles to detect \vcs (\Cref{sec:guided-search}). 

\subsection{Stitching Causal Relationships}
\label{sec:chain-forming}

\tool stitches causal relationships 
when the resulting fault in one causal relationship 
($f_2$ in $f_1 \rightarrow f_2$ ) 
matches the starting fault in another ($f_2$ in $f_2 \rightarrow f_3$). 
In particular, 
based on the type of the fault ($f_2$) used in stitching, 
\tool can form four types of connections: 
\begin{enumerate}[leftmargin=3ex, noitemsep, topsep=0pt]
\item $f_1 \longEArrow{D/I} f_2 \longAnyInterferenceArrow{I} f_3$: 
The matching fault $f_2$ is an exception or negation. Injecting $f_1$ has an execution trace interference $f_2$ on the system. $f_1$ and $f_3$ can be of any type.

\item $f_1 \longSPlusArrow{D/I} f_2 \longAnyInterferenceArrow{D} f_3$:
The matching fault $f_2$ is a delay on a loop. Injecting $f_1$ has an iteration count interference $f_2$ on the system. $f_1$ and $f_3$ can be of any type. 
For delays in nested loops performing batch processing, there are two variants: 

\begin{enumerate}[leftmargin=3ex]
\item $f_1 \longSPlusArrow{D/I}$ \CFGBackEdge{$f_2$}{$f_2'$}$\longAnyInterferenceArrow{D} f_3$:
The delay $f_2$ propagates to its parent loop ($f_2'$), which is the injected fault of the next injection and causes $f_3$.

\item $f_1 \longSPlusArrow{D/I}$ \CFGBackEdge{$f_2$}{} \CFGEdge{$f_2'$}{$f_2''$}$\longAnyInterferenceArrow{D} f_3$:
The delay $f_2$ propagates to its sibling loop ($f_2''$), which is the injected fault of the next injection and causes $f_3$.
\end{enumerate}
\end{enumerate}

\subsection{Local Compatibility Check}
\label{sec:chain-forming:approx-check}

\begin{figure}
\centering
\footnotesize
\begin{minted}[
  xleftmargin=2ex, xrightmargin=1ex, frame=lines, fontsize=\footnotesize, linenos, 
  autogobble, numbersep=1ex, framesep=1mm, breaklines]
{java}
public ReplicaHandler createTmp(...) {
  do {
    if (current == lastFoundReplica) { ... };
    // ...
    // Fault F2 (Injected)
    if ((current.genStamp() >= b.genStamp()) || !isTransfer)
      // Fault F2 (Triggered by F1 injection)
      throw new ReplicaAlreadyExistsException(...);
  } while (true);
}

class BlockReceiver {
  BlockReceiver(...){
    data.createTmp(...);
  }
}
\end{minted}
% \vspace{-3mm}
\caption{Pseudo-code demonstrating state compatibility. The fault is located inside a loop of \texttt{createTmp()}, which is invoked from \texttt{BlockReceiver()}. Code simplified from \texttt{FsDatasetImpl} in HDFS.}
% \vspace{-2mm}
\label{fig:state-compatibility}
\end{figure}


\noindent
When connecting causal relationships of faults obtained from different tests, \tool performs a compatibility check 
between the tests involved. 
Theoretically, 
to check for 
incompatible propagation activation conditions across tests (and skip stitching), 
a symbolic constraint collection 
(i.e., path condition collection used in symbolic execution techniques~\cite{cadar2008klee,chipounov2012s2e})
should be applied during the fault propagation: 
collect the path conditions $c_1$ and $c_2$
during the propagation of  $f_1 \rightarrow f_2$ in test $t_1$
and $f_2 \rightarrow f_3$ in test $t_2$ respectively, 
and check the satisfiability of their conjunction 
$c_1 \wedge c_2$.

To minimize the overhead, 
\tool approximates this analysis 
by checking whether 1) the \textbf{local} execution trace associated with the fault 
used in stitching ($f_2$) and 2) the call stack 
match with their counterparts in the other test.
In practice, 
we find \tool's strategies successfully eliminates 
most incompatible workloads. 

To facilitate the compatibility check, 
in 
injection runs, 
\tool records call stack and execution trace as follows: 
\begin{enumerate}[wide, labelwidth=!, labelindent=\parindent, noitemsep, topsep=0pt]

\item
\textit{Execution trace recording}. 
\tool records the branch statements and the evaluation results of the conditions 
\textbf{locally} in the fault's enclosing loop or function 
as the path conditions. 
For example, in \Cref{fig:injection-point-exp}, the three \texttt{if}s (line 3, 7, and 15) serve as monitor points for such branch statements. 

\item
\textit{Call stack recording}. 
An error occurred in the same function but at different call sites 
may represent errors for different types of requests, 
which is invalid when stitched together. 
To address this, \tool records the closest two levels of call stack for each loop iteration, exception, and boolean return value (excluding the enclosing function itself). 
This method resembles the  2-call-site sensitivity in pointer analysis, a common method balancing speed and accuracy~\cite{jeon2022return,li2020principled, tan2017efficient}. 
Note that 
the scope of 
execution trace 
and call stack 
used in the comparison 
can both be adjusted. 

\end{enumerate}


\Cref{fig:state-compatibility} shows 
an example for the compatibility check. 
The fault $f_2$ (exception on line 8) used in stitching 
is inside a loop of \texttt{createTmp()}.  
\tool checks the following two things before stitching $f_1 \rightarrow f_2$ and $f_2 \rightarrow f_3$ together: 
1) (call stack) \texttt{create\-Tmp()} must be both invoked from \texttt{Block\-Re\-ceiver()}; 
2) (execution trace) 
the traces of line 3-5 
in the 
fault-happening 
iteration 
of the fault-enclosing loop match between tests. 
If the fault is not within a loop, 
we use the trace of the enclosing function (e.g., \texttt{Block\-Re\-ceiver()}) instead. 
Because delay injection is injected at the beginning of all iterations of a loop, 
\tool conservatively checks 
for matching traces in any loop iterations between tests. 

